% =============================================================================
% chapter_conclusion.tex — General Conclusion
% =============================================================================

\chapter*{General Conclusion}
\addcontentsline{toc}{chapter}{General Conclusion}
\markboth{General Conclusion}{General Conclusion}
\thispagestyle{fancy}

This project set out to answer a question that occupies data engineering teams
across many industries: can an automated system translate enterprise SAS code to
Python at a quality level that is useful in production --- not merely syntactically
plausible, but semantically verifiable? The answer, as demonstrated by the
\codara{} system, is: yes, for the majority of real-world SAS constructs, and
with formally grounded confidence signals that tell the engineer exactly when to
trust the output and when to review it manually.

\textbf{What was built.} \codara{} is a multi-agent SAS-to-Python conversion
accelerator comprising an eight-node LangGraph orchestration pipeline, three
complementary RAG retrieval tiers, a 330-pair LanceDB knowledge base, Z3 SMT
formal verification with CEGAR repair, the CDAIS adversarial input synthesis
framework, the SemantiCheck four-layer semantic scoring system, a FastAPI
REST backend, and a React web dashboard --- deployed as a containerised service
on Azure Container Apps with a six-job GitHub Actions CI/CD pipeline.

\textbf{What was demonstrated.} The evaluation on the 50-file, 721-block
gold-standard corpus (45 with hand-verified translations) confirmed six of the
eight technical targets: 72.4\% translation success rate (target: 70\%),
96.8\% syntax-valid output (target: 95\%), RAPTOR retrieval advantage of 12
percentage points on MODERATE/HIGH partitions (target: $\geq$10 pp), ECE of
0.061 (target: $<$0.08), Z3 SMT provability of 41\% (target: $\geq$35\%), and
streaming performance of 1.34 seconds / 61 MB for a 10,000-line file
(targets: $<$2 s / $<$100 MB). The 10-block torture test achieved 10/10
execution success with 3 Z3 formal proofs and 5 blocks at VERIFIED or
LIKELY\_CORRECT SemantiCheck status.

Two targets were not met: boundary detection F1 of 79.3\% fell short of the
90\% target, and hard-tier translation success at 53.7\% reflects the
unresolved macro meta-programming challenge. Both trace to the same root cause:
the SAS macro language requires a full pre-expansion pass before structural
analysis, which is the primary item on the short-term improvement roadmap.

\textbf{What was novel.} The five scientific contributions of this work are:
(1)~three-tier RAG with RAPTOR hierarchical clustering, demonstrated to
outperform flat-index retrieval by 12 pp on complex SAS constructs;
(2)~Z3 SMT formal verification with CEGAR repair applied to LLM-generated
Python translations, catching 43 real semantic errors undetectable by execution
testing; (3)~CDAIS adversarial input synthesis with Z3-driven minimum-witness
construction and formal coverage certificates; (4)~SemantiCheck, a four-layer
composite semantic verification framework combining formal proofs, adversarial
tests, STG contract checking, and LLM-as-oracle simulation; and (5)~LLM-as-SAS-oracle,
a technique for behavioural testing of SAS translations without a SAS licence.

\textbf{Looking forward.} The fifteen-week sprint produced a production-ready
v3.1.0 system and a clear research agenda. The short-term priority is a SAS
macro pre-processor that enables boundary detection before rather than after
macro call sites. The medium-term research directions --- a QLoRA fine-tuned
local model, a full SAS grammar parser for STG extraction, and CDAIS extension
to eight error classes --- build on the framework established here to improve
both quality and coverage. The long-term vision is a domain-agnostic translation
accelerator that applies the same principles (three-tier RAG, adversarial
synthesis, composite semantic scoring) to other legacy-to-modern migration
challenges: COBOL-to-Java, Fortran-to-Python, legacy PL/SQL to cloud SQL
dialects.

The most important lesson of this project is neither technical nor architectural.
It is methodological: the combination of Design Science Research and Agile
development proved more productive than either would have been alone. The DSR
framework forced us to state claims precisely enough to test them, and the Agile
sprints ensured that every two weeks, there was working software that could
generate real evidence. That discipline --- claim, build, measure, revise ---
is what distinguishes an engineering contribution from an engineering exercise.
